\documentclass[a4paper, 11pt]{article}
\usepackage[T1]{fontenc}
\usepackage{mathpazo}
\usepackage[hmargin=1in,vmargin=1in]{geometry}
\usepackage{amsmath,amssymb,amsthm}
\usepackage{graphicx}
\graphicspath{{manuscript/}}
\usepackage{longtable}
\usepackage{booktabs}
\usepackage{array}
\usepackage{rotating}
\usepackage{setspace}
\usepackage{subcaption}
\usepackage{xcolor}
\usepackage{authblk}
\usepackage[hyphens]{url}
\usepackage{microtype}
\usepackage{pdflscape}
\usepackage{etoolbox}
\usepackage[natbibapa]{apacite}
\usepackage[hidelinks]{hyperref}
\hyphenpenalty=10000
\setlength{\emergencystretch}{3em}
\renewcommand{\topfraction}{0.85}
\renewcommand{\textfraction}{0.10}
\renewcommand{\floatpagefraction}{0.80}
\exhyphenpenalty=10000
\interfootnotelinepenalty=10000
\widowpenalty=10000
\clubpenalty=10000
\AtBeginEnvironment{thebibliography}{%
  \interlinepenalty=10000
}

\date{October 2026}

\singlespacing

\author[1]{Timothy B. Gravelle}
\author[2]{Stuart Bagnall}
\author[2]{Antonios Katris}
\author[2]{Frankie Watson-Brown}
\author[3]{Jason Reifler}
\author[2]{Thomas J. Scotto}
\affil[1]{Wilfrid Laurier University}
\affil[2]{University of Strathclyde}
\affil[3]{University of Southampton}

\title{Nuclear Energy in the Neighbourhood?\\Examining Public Support for Small Modular Reactors in the United Kingdom}

\begin{document}

\pagenumbering{roman}
\thispagestyle{empty}
\clearpage

\maketitle

\begin{abstract}
\normalsize
\noindent The United Kingdom faces growing demands on its electricity grid, and electricity prices have become a major cost of living issue. At the same time, the UK has made a commitment to transition away from carbon-emitting sources of energy, and to build small modular reactors (SMRs). How SMRs are described and where they are sited can prompt considerations that lead people to support or oppose them. Siting matters in two senses: where SMRs are said to be built and how close people actually live to nuclear sites, as proximity may induce opposition by heightening perceptions of risk or support through expected economic benefits. This paper explores these issues using data from an original survey-based experiment (\textit{n} = 1,276). The experiment described SMRs as ``reactors'' or ``power stations'' developed ``in the United Kingdom'' or ``in your local area''. Support is favourable but not overly so. Calling SMRs reactors rather than power stations shifts support only slightly, while describing development as local lowers it modestly, consistent with ``NIMBY'' accounts. There are also regional dynamics at work, with lower support in Scotland compared to England, especially when the question relates to developing modular reactors in the local area. Lower support in Scotland is concentrated among supporters of independence. Support is also lower in those parts of the UK near existing nuclear sites as opposed to those further away.\\

\textit{Keywords:} Public opinion, nuclear energy, small modular reactors, framing, NIMBY, devolution, United Kingdom

\end{abstract}

\end{document}